\section{Limitations and open tests}
\label{sec:limitations}

Appendix~\ref{sec:limitations-full} expands each point. \textbf{Data:} WildChat is consumer chat, not API traffic; hashed IPs only approximate accounts; activity thresholds exclude low-volume accounts, where a patient attacker would hide; and colluders are 14--50\% of a window, far above production shares. \textbf{Representation:} attacker and defender share one family of sentence embeddings (E5-small-v2 is initialised from MiniLM); lexical, style, timing and victim-model features are untested. \textbf{Statistics:} conclusions rest on 17--24 test weeks; development results reflect adaptive reuse of the same weeks, which one frozen evaluation only partly corrects; predictions were locked in notebooks but carry no external timestamp. \textbf{Campaign construction:} colluder accounts were assembled from draws across a pool rather than from whole real user-weeks, which confounds the natural and covert results in an unknown direction. \textbf{Capability:} the frontier rests on one 360M student, one benchmark, answers from datasets rather than a live teacher, and two seeds; the imitating writer was a 1.5B model.

As future work, the experiments that would most change our conclusions are an anytime-valid sequential test on the hollowed score, colluders built from whole real user-weeks, a relay simulation, a defence against calibration poisoning, distillation from a live teacher with more seeds, students and tasks, and, above all, the benign-coordination census on a real provider's logs.
